\documentclass[11pt]{article}
\usepackage{ochamath}
\subjectcolor{2F5DA8}
\setsheet{線形代数}{ベクトルと内積　演習}{https://math.ochanote.com/linear-algebra/vectors/}
\namelinetrue
\begin{document}
\makesheethead{問題}
この演習では実ベクトルの標準内積とユークリッドノルムを使う。

\problem[成分の計算と長さ]
$\boldsymbol a=\begin{pmatrix}1\\2\end{pmatrix}$、$\boldsymbol b=\begin{pmatrix}3\\-1\end{pmatrix}$ とする。
\begin{enumerate}[label=(\arabic*)]
\item $\boldsymbol a+\boldsymbol b$ と $2\boldsymbol a-\boldsymbol b$ を求めよ。
\item $\boldsymbol a$ の長さと、$\boldsymbol a$ と同じ向きの単位ベクトルを求めよ。
\end{enumerate}
\workspace{45mm}

\problem[内積の意味]
$\boldsymbol u=\begin{pmatrix}1\\1\end{pmatrix}$、$\boldsymbol v=\begin{pmatrix}2\\0\end{pmatrix}$、$\boldsymbol w=\begin{pmatrix}1\\-1\end{pmatrix}$ とする。
\begin{enumerate}[label=(\arabic*)]
\item $\boldsymbol u\cdot\boldsymbol v$ と、$\boldsymbol u$ と $\boldsymbol v$ のなす角を求めよ。
\item $\boldsymbol u\cdot\boldsymbol w$ を求め、2つのベクトルの向きの関係を説明せよ。
\end{enumerate}
\workspace{45mm}

\newpage
\problem[零ベクトルへの注意]
次の主張について、正しいかどうかを理由とともに答えよ。
\begin{enumerate}[label=(\arabic*)]
\item 任意の実ベクトル $\boldsymbol a$ について $\boldsymbol a\cdot\boldsymbol 0=0$ である。
\item 零ベクトルを長さで割ると、単位ベクトルが得られる。
\item $\boldsymbol a\cdot\boldsymbol b=0$ なら、どんな場合でも2つのベクトルのなす角は $90^\circ$ といえる。
\end{enumerate}
\workspace{50mm}

\problem[一定の力がする仕事]
物体に一定の力 $\boldsymbol F=\begin{pmatrix}4\\-3\end{pmatrix}\ \mathrm N$ が加わり、
物体の変位が $\boldsymbol d=\begin{pmatrix}2\\1\end{pmatrix}\ \mathrm m$ であった。
\begin{enumerate}[label=(\arabic*)]
\item この力がする仕事 $W$ を求めよ。
\item $\|\boldsymbol F\|\,\|\boldsymbol d\|$ を計算し、その値が仕事と異なる理由を説明せよ。
\item 同じ力のもとで、変位が $-\boldsymbol d$ だった場合の仕事を求めよ。
\end{enumerate}
\workspace{65mm}
\end{document}
